[V.] Art. $[7.]$ Here, as in Article $1$ of the preceding fragment [IV], $\mathrm{Dec.}\frac mM$ denotes the value of
\[
\Sigma(-1)^X\Sigma(-1)^Y\mathrm{Int.}\,p
\]
where the summation extends over all integers $X$ and $Y$ for which the corresponding $\xi$ and $\eta$ lie within the bounds $0$ and $\tfrac12$.

The text states that the formulas for the decident in Articles $7$ and $15$ were found in two distinct
\SourceLeaf{388}
ways. To keep these explanations from becoming too long, however, they are derived here from a common source.

For any fixed integer value of $X$, let $Y^{*}$ be the smaller and $Y^{**}$ the larger of the two values of $Y$ corresponding to the boundary values of $\xi$, $\eta$. Let the values of $p$ corresponding to $Y^{*}$ and $Y^{**}$ be $p^{*}$ and $p^{**}$. Like $Y^{*}$ and $Y^{**}$, they cannot be equal to one another, since the sum $\Sigma$ does not extend over the boundary values of $\xi$ and $\eta$.

Sum over all values of $Y$ occurring for the same $X$, in the same way as in the preceding two papers [III] and [IV]. Put $[\tfrac12Y''-\tfrac12]-[\tfrac12Y'-\tfrac12]$ for the number of odd integers between $Y'$ and $Y''$, and add the intensors referring to the bounds $\xi=0$ and $=\tfrac12$ twice, with opposite signs. This gives for $\Sigma(-1)^Y\mathrm{Int.}\,p$ the expression consisting of seven parts
\[
\begin{gathered}
-\Sigma[-\mathrm{Int.}(p-\mu\omega i)+\mathrm{Int.}(p+\mu\omega i)]\\
\text{taking all }p\text{ for which }[Y]\text{ even,}\\
x\text{ or }y\text{ integer, incl. }\xi=0\text{ and }\tfrac12,\quad
\text{excl. }\eta=0\text{ and }\tfrac12
\end{gathered}
\]
\[
\left.\begin{aligned}
&-\mathrm{Int.}(p^{*}-\mu\omega i)\text{ if }[Y^{*}]\text{ even}\\
&+\mathrm{Int.}(p^{**}+\mu\omega i)\text{ if }[Y^{**}]\text{ even}
\end{aligned}\right\}\ \xi=0\text{ or }\tfrac12,\quad0<\eta<\tfrac12
\]
\[
\left.\begin{aligned}
&-\mathrm{Int.}(p^{*}+\mu\omega i)\text{ if }[Y^{*}]\text{ even}\\
&+\mathrm{Int.}(p^{**}-\mu\omega i)\text{ if }[Y^{**}]\text{ even}
\end{aligned}\right\}\ 0<\xi<\tfrac12,\quad\eta=0\text{ or }\tfrac12
\]
\[
\left.\begin{aligned}
&-\mathrm{Int.}(p^{*}+\mu\omega i)\text{ if }[Y^{*}]\text{ even}\\
&+\mathrm{Int.}(p^{**}-\mu\omega i)\text{ if }[Y^{**}-\omega]\text{ even}
\end{aligned}\right\}\ \xi=0\text{ or }\tfrac12,\quad\eta=0\text{ or }\tfrac12
\]
which, multiplied by $(-1)^X$ and summed over all integer values of $X$, gives the decident $\frac mM$.

In this way, the first part of the expression gives, from the $\varepsilon$ formed according to Article $2$, Rule I,
\[
\begin{aligned}
&\Sigma\varepsilon\text{ where }X\text{ integer, }[Y]\text{ even}\\
&-\Sigma\varepsilon\text{ where also }y\text{ integer, }[x]\text{ even}
\end{aligned}
\]
For the following parts one may put
\[
\left.\begin{aligned}
&-(B)\mathrm{Int.}(p-B\mu\omega i)\text{ if }\xi=0\\
&+(B)\mathrm{Int.}(p+B\mu\omega i)\text{ if }\xi=\tfrac12
\end{aligned}\right\}\ X\text{ integer, }[Y]\text{ even, }0<\eta<\tfrac12
\]
\[
\left.\begin{aligned}
&-(A)\mathrm{Int.}(p+A\mu\omega i)\text{ if }\eta=0\\
&+(A)\mathrm{Int.}(p-A\mu\omega i)\text{ if }\eta=\tfrac12
\end{aligned}\right\}\ X\text{ integer, }[Y]\text{ even, }0<\xi<\tfrac12
\]
\[
\begin{gathered}
-\tfrac12[(A')+(B')]\mathrm{Int.}(p+A'\mu\omega i)\\
\text{if }X\text{ integer, }[Y+A'\omega]\text{ even, }\xi\text{ and }\eta=0\text{ or }\tfrac12
\end{gathered}
\]
where $A'=+A$ or $-A$ when $\eta=0$ or $\tfrac12$, and $B'=+B$ or $-B$ when $\xi=0$ or $\tfrac12$; here, for example, the notation $(A):+1$ or $-1$ is used according as $A$ is positive or negative.

Multiply by $(-1)^X$ and sum over $X$, making the following substitutions in these expressions, namely in the
\[
\begin{aligned}
\text{first}\quad&p-B\mu\omega i,\ P-BD\omega i,\ \pi-AB\omega i-BB\omega,\ X,\ [Y]\\
&\text{respectively into }ip,\ iP,\ i\pi,\ -Y,\ [X]\\[3pt]
\text{second}\quad&p+B\mu\omega i,\ P+BD\omega i,\ \pi+AB\omega i+BB\omega,\ X,\ [Y]\\
&\dotstwo\quad p,\ P,\ \pi\quad X,\ [Y]\\[3pt]
\text{third}\quad&p+A\mu\omega i,\ P+AD\omega i,\ \pi+AA\omega i+AB\omega,\ X,\ [Y]\\
&\dotstwo\quad p,\ P,\ \pi\quad X,\ [Y]\\[3pt]
\text{fourth}\quad&p-A\mu\omega i,\ P-AD\omega i,\ \pi-AA\omega i-AB\omega,\ X,\ [Y]\\
&\dotstwo\text{ into }im-ip,\ iM-iP,\ i-i\pi,\ Y-B,\ A-1-[X]
\end{aligned}
\]

\SourceLeaf{389}
and denote the result from the fifth expression by $Q_1$. One then obtains
\[
\begin{aligned}
&-\Sigma(-1)^Y(B)\mathrm{Int.}\,ip,\\
&\qquad\text{where }Y\text{ integer, }[X]\text{ even, }\eta=\omega,\ 0<\xi<\tfrac12\\
&+\Sigma(-1)^X(B)\mathrm{Int.}\,p,\\
&\qquad\text{where }X\text{ integer, }[Y]\text{ even, }\xi=\tfrac12+\omega,\ 0<\eta<\tfrac12\\
&-\Sigma(-1)^X(A)\mathrm{Int.}\,p,\\
&\qquad\text{where }X\text{ integer, }[Y]\text{ even, }\eta=\omega,\ 0<\xi<\tfrac12\\
&+\Sigma(-1)^Y(A)\mathrm{Int.}(im-ip),\\
&\qquad\text{where }Y\text{ integer, }[X]\text{ even, }\xi=\tfrac12+\omega,\ 0<\eta<\tfrac12\\
&+Q_1
\end{aligned}
\]

Examining the individual cases shows that, under the assumption $M\equiv1\ (\text{mod. }2+2i)$,
\[
\begin{aligned}
Q_1={}&-\mathrm{Int.}\,\mu\omega i\text{ if }M\text{ in quadrant }1\text{}\\
&-\mathrm{Int.}(\tfrac12mi+\mu\omega i),\quad\text{if }M\text{ in quadrant }2\text{ and }\frac{M-1}{2}\text{ even}\\
&+\mathrm{Int.}(\tfrac12mi-\mu\omega i),\quad\text{if }M\text{ in quadrant }4\text{ and }\frac{M-1}{2}\text{ even}
\end{aligned}
\]
where a complex number is called even or odd according as it is divisible by $2$ or not.

Thus, using the $\varepsilon$ specified in Rules II and IV, one obtains
\[
\begin{aligned}
\mathrm{Dec.}\frac mM={}&\mathrm{I},& &\Sigma\varepsilon,\ \text{where }X\text{ integer, }[Y]\text{ even}\\
&&&-4\Sigma\varepsilon,\ \text{where also }y\text{ integer, }[x]\text{ even}\\
&\mathrm{II},&&-\Sigma\varepsilon\mathrm{Int.}\,ip,\ \text{where }Y\text{ integer, }[X]\text{ even}\\
&\mathrm{IV},&&+\Sigma\varepsilon\mathrm{Int.}\,p,\ \text{where }X\text{ integer, }[Y]\text{ even}\\
&\mathrm{II},&&+\Sigma\varepsilon\mathrm{Int.}\,p,\ \text{where }X\text{ integer, }[Y]\text{ even}\\
&\mathrm{IV},&&+\Sigma\varepsilon\mathrm{Int.}(im-ip),\ \text{where }Y\text{ integer, }[X]\text{ even}\\
&&&+Q_1
\end{aligned}
\]

The expression for the decident takes another form if one first sums over $X$, expressing the number of odd integers between $X'$ and $X''$ by $[\tfrac12X''+\tfrac12]-[\tfrac12X'+\tfrac12]$, namely
\[
\begin{aligned}
\mathrm{Dec.}\frac mM={}&\mathrm{I},& &\Sigma\varepsilon,\ \text{where }Y\text{ integer, }[X]\text{ odd}\\
&&&-4\Sigma\varepsilon,\ \text{where also }y\text{ integer, }[x]\text{ even}\\
&\mathrm{II},&&-\Sigma\varepsilon\mathrm{Int.}\,ip,\ \text{where }X\text{ integer, }[Y]\text{ even}\\
&\mathrm{IV},&&+\Sigma\varepsilon\mathrm{Int.}\,p,\ \text{where }Y\text{ integer, }[X]\text{ odd}\\
&\mathrm{II},&&+\Sigma\varepsilon\mathrm{Int.}\,p,\ \text{where }Y\text{ integer, }[X]\text{ odd}\\
&\mathrm{IV},&&+\Sigma\varepsilon\mathrm{Int.}(im-ip),\ \text{where }X\text{ integer, }[Y]\text{ odd}\\
&&&+Q_2
\end{aligned}
\]
\[
\begin{aligned}
Q_2={}&-\mathrm{Int.}(-\mu\omega),\quad\text{if }M\text{ in quadrant }2\text{}\\
&+\mathrm{Int.}(\tfrac12m-\mu\omega),\quad\text{if }M\text{ in quadrant }1\text{ and }\frac{M-1}{2}\text{ odd}\\
&-\mathrm{Int.}(\tfrac12m+\mu\omega),\quad\text{if }M\text{ in quadrant }3\text{ and }\frac{M-1}{2}\text{ odd}
\end{aligned}
\]

\SourceLeaf{390}
If one sums over $Y$ first, but uses the second method just applied for determining the number of odd integers between two values, one obtains
\[
\begin{aligned}
\mathrm{Dec.}\frac mM={}&\mathrm{I},& &\Sigma\varepsilon,\ \text{where }X\text{ integer, }[Y]\text{ odd}\\
&&&-4\Sigma\varepsilon,\ \text{where also }y\text{ integer, }[x]\text{ even}\\
&\mathrm{II},&&-\Sigma\varepsilon\mathrm{Int.}\,ip,\ \text{where }Y\text{ integer, }[X]\text{ odd}\\
&\mathrm{IV},&&+\Sigma\varepsilon\mathrm{Int.}\,p,\ \text{where }X\text{ integer, }[Y]\text{ odd}\\
&\mathrm{II},&&+\Sigma\varepsilon\mathrm{Int.}\,p,\ \text{where }X\text{ integer, }[Y]\text{ odd}\\
&\mathrm{IV},&&+\Sigma\varepsilon\mathrm{Int.}(im-ip),\ \text{where }Y\text{ integer, }[X]\text{ odd}\\
&&&+Q_3
\end{aligned}
\]
\[
\begin{aligned}
Q_3={}&-\mathrm{Int.}(-\mu\omega i),\quad\text{if }M\text{ in quadrant }3\text{}\\
&-\mathrm{Int.}(\tfrac12mi+\mu\omega i),\quad\text{if }M\text{ in quadrant }2\text{ and }\frac{M-1}{2}\text{ odd}\\
&+\mathrm{Int.}(\tfrac12mi-\mu\omega i),\quad\text{if }M\text{ in quadrant }4\text{ and }\frac{M-1}{2}\text{ odd}
\end{aligned}
\]

If one also sums over $X$ first and uses the first method for expressing the number of odd integers between two values, one obtains the form for the decident given in [V], Article $15$, where the quantity $q$ can also be defined by the following equation:
\[
\begin{aligned}
q={}&-\mathrm{Int.}\,\mu\omega,\quad\text{if }M\text{ in quadrant }4\text{}\\
&+\mathrm{Int.}(\tfrac12m-\mu\omega),\quad\text{if }M\text{ in quadrant }1\text{ and }\frac{M-1}{2}\text{ even}\\
&-\mathrm{Int.}(\tfrac12m+\mu\omega),\quad\text{if }M\text{ in quadrant }3\text{ and }\frac{M-1}{2}\text{ even}
\end{aligned}
\]

Combining these four expressions for the decident gives the result stated in [V], Article $7$, since $\mathrm{Int.}\,ip-\mathrm{Int.}\,p$ equals $3$ when $[x]$ is even and $[y]$ odd, and otherwise equals $1$; furthermore, $\mathrm{Int.}(im-ip)+\mathrm{Int.}\,p$ equals $0$ when $[x]$ and $[y]$ are even, and equals $4$ in all other cases.

[V.] Art. $[7.]$ In the first column, the first table for the example gives the values of $\frac{37\pdot P}{M}$ or $37(\xi+\eta i)$ corresponding to each integer $P$, when $0<\xi<\tfrac12$, $0<\eta<\tfrac12$; in the second, $\frac{37\pdot Pm}{M}$ or $37\pdot p$; in the third, the greatest integer contained in $p$; and in the fourth, $\pm\mathrm{Int.}\,p$, where the upper sign applies when $P$ is divisible by $1+i$, and the lower when $P$ is not divisible by $1+i$.

[V.] Art. $[9.]$. $[12.]$ The improved notation for the $\theta$ is indicated only for the second and third classes, Articles $9$. $10$, but is also to be extended to the first and fourth, Articles $8$. $11$. Accordingly, a $\theta^\lambda=T+Ui$ has the index $\lambda,=0,1,2$ or $3$ for which the quantities $\varphi^0$ and $\Phi^0$, determined by the equations

\SourceLeaf{391}
\[
\begin{aligned}
i^\lambda M&=\mathfrak A+\mathfrak B i,\quad i^{-\lambda}\mu=\rho+\sigma i\\
\sigma\varphi^0&=-\text{Coeff. Imag. }\theta^\lambda(a-bi)=+bT-aU\\
\sigma\Phi^0&=+\text{Coeff. Imag. }\theta^\lambda(\mathfrak A-\mathfrak B i)=-\mathfrak B T+\mathfrak A U
\end{aligned}
\]
lie between $0$ and $\tfrac12$.

To prove, following the indications in Article $9$ $(3)$, that if $T$, $U$ are two real integers satisfying the conditions just stated, then $T+Ui$ also belongs to the collection $\theta^\lambda$ determined in one of the four combinations, Articles $8$. $11$, let $\varphi'$, $\Phi'$ denote the complex integers for which the equation
\[
T+Ui=\varphi'i^\lambda M+\Phi'm
\]
holds, and for which one of the four quantities $\pm\frac{\varphi'-\varphi^0}{m}$, $\pm i\frac{\varphi'-\varphi^0}{m}$ has its real part and the coefficient of its imaginary part between $0$ and $\tfrac12$ (\textit{Theoria residuorum biquadr. artt.} $45,46$). As is readily seen by examining the sixteen individual cases in the four combinations, the quantity in question is then $\frac pm$, and the corresponding quantity among $\pm\frac{\Phi'-\Phi^0}{m}$, $\pm i\frac{\Phi'-\Phi^0}{M}$ is $\frac PM$, since $\frac{\Phi'-\Phi^0}{M}=-\frac{\varphi'-\varphi^0}{m}i^\lambda$.

This way of representing the quantities $\frac pm$ or $\frac PM$ also shows that the sum $\mathrm I,\ \Sigma\varepsilon$, over all, equals $\Theta+\Theta'+\Theta''+\Theta'''$, where $\Theta^\lambda$ denotes the sum of the $\varepsilon$ which, for each lattice point $\theta$ inside the parallelogram $0,\ \tfrac12m,\ \tfrac12m+\tfrac12i^\lambda M,\ \tfrac12i^\lambda M,$ are to be put
\begin{itemize}
\item[] $=+1$ if $\theta$ is divisible by $1+i$ and $\text{Coeff. Imag. }\mu i^{-\lambda}$ is positive, or if neither condition holds; but
\item[] $=-1$ if only one holds.
\end{itemize}

[V.] Art. $13$. The quantity $\Theta^\lambda$ can be determined either by the four different ways of summing used above for $\mathrm{Dec.}\frac mM$, or, what is essentially the same, by the methods indicated in [II], Article $11$. These are to be based on the four constructions obtained by joining the points whose $\theta$ is a multiple of $1+i$ to the points $\theta+1$, $\theta+i$, $\theta-1$ and $\theta-i$, respectively.

On the boundary of the parallelogram just mentioned, let a $\theta$ correspond to every point for which the real or imaginary part of $\theta$ is an integer. Denote by $\theta^0$ the nearest integer divisible by $1+i$ to $\theta$, and by $l$ the displacement from a point on the straight boundary segment containing $\theta$ to any following point on the same straight line; thus, for example, for this parallelogram, the quantities $m,\ Mi^\lambda,\ -m,\ -Mi^\lambda$, in order. If one puts
\[
\varepsilon=\pm1\text{ with the sign of the imaginary part of }\frac l{\theta-\theta^0}
\]
then combining the four results obtained in either way gives $4\Theta^\lambda=-\Sigma\varepsilon$.

One avoids determining the $\theta$ and $\varepsilon$ at the vertices separately by replacing this parallelogram with another whose boundaries are infinitely near those of the former
\SourceLeaf{392}
and which does not include the two points $0$ and $\tfrac12m+\tfrac12i^\lambda M$. The boundary of such a parallelogram is obtained by placing it on the positive sides of the lines
\[
\begin{gathered}
0\dotsfour\tfrac12m,\quad\tfrac12m+\tfrac12i^\lambda M\dotsfour\tfrac12m,\\
\tfrac12m+\tfrac12i^\lambda M\dotsfour\tfrac12i^\lambda M,\quad 0\dotsfour\tfrac12i^\lambda M
\end{gathered}
\]
Assign the infinitely small positive quantities $\omega_1,\omega_2,\omega_3,\omega_4$ to the four resulting straight lines, in order. For the points $\theta$ on them one may then put
\[
\begin{aligned}
\theta&=p=m(\xi+\omega_1i),\\
-\theta i^\lambda+\tfrac12mi^\lambda+\tfrac12M&=P=M(\xi+\omega_2i),\\
-\theta+\tfrac12m+\tfrac12Mi^\lambda&=p=m(\xi+\omega_3i),\\
\theta i^{-\lambda}&=P=M(\xi+\omega_4i)\quad\text{if }\lambda\text{ even}
\end{aligned}
\]
\[
\begin{aligned}
\theta&=p=m(\xi+\omega_1i),\\
\theta i^{1-\lambda}-\tfrac12mi^{1-\lambda}+\tfrac12M&=P=M(\tfrac12+\omega_2+\eta i),\\
\theta i+\tfrac12m-\tfrac12Mi^{1+\lambda}&=p=m(\tfrac12+\omega_3+\eta i),\\
\theta i^{-\lambda}&=P=M(\xi+\omega_4i)\quad\text{if }\lambda\text{ odd}
\end{aligned}
\]
where, in part, $\xi$ and $\eta$ also extend beyond the range of real values from $0$ to $\tfrac12$ by infinitely small quantities, to close the figure.

Let $G$, $g$, $H$, $h$ denote the sums of the $\varepsilon$ formed according to Rules II, III, IV, V, respectively (in Articles $3$ to $6$). Let $G'$ or $G_{\prime}$, and $g'$ or $g_{\prime}$, include the $\varepsilon$ arising for the infinitely small values of $\xi$ that may occur in the second parallelogram, but otherwise formed according to Rules II and III, respectively. Further, let $G''$ or $G_{\prime\prime}$, and $g''$ or $g_{\prime\prime}$, similarly refer to the same rules, but to the infinitely small values of $\tfrac12-\xi$; and finally let $H',h_{\prime},H'',h_{\prime\prime}$ refer to Rules IV, V, IV, V and the infinitely small values $\tfrac12-\eta,\tfrac12-\eta,\eta,\eta$, respectively. Then
\[
\begin{aligned}
4\Theta^\lambda={}&-(g+g'+g'')-i^\lambda(G+G'+G'')\\
&+i^\lambda(g+g_{\prime}+g_{\prime\prime})+(G+G_{\prime}+G_{\prime\prime})\quad\text{if }\lambda\text{ even}\\
4\Theta^\lambda={}&-(g+g'+g'')-i^{1-\lambda}(H+H'+H'')\\
&-i^{\lambda-1}(h+h_{\prime}+h_{\prime\prime})+(G+G_{\prime}+G_{\prime\prime})\quad\text{if }\lambda\text{ odd}
\end{aligned}
\]

For the vertex $\theta$ of the parallelogram nearest to $0$, let $\xi_1$ denote the corresponding value of the $\xi$ of the first side, and $\xi_4$ the corresponding value of the $\xi$ of the fourth side, so that
\[
\theta=m(\xi_1+\omega_1i)=i^\lambda M(\xi_4+\omega_4i)
\]
Then the $\xi$ corresponding to a point $p$ on the first side or its extension with real part equal to $0$ is given by
\[
(\mathrm{Real.}\,p=0),\quad\xi-\xi_1=\sigma a\mathfrak A\omega_1-\sigma\omega_4
\]
where the positive factors of the infinitely small positive quantities have been replaced by unity, and $\sigma$, $\mathfrak A$ denote the real quantities determined by
\[
\rho+\sigma i=i^{-\lambda}(\alpha+\beta i),\quad\mathfrak A+\mathfrak B i=i^\lambda(A+Bi)
\]
This point $p$ lies on the first side itself if $\xi-\xi_1$ is positive; hence, taking $\omega_1$ infinitely small compared with $\omega_4$, if $\sigma$ is negative. The point $p^0$ nearest to $p$ whose representing number is divisible by $1+i$ is $0$; thus
\SourceLeaf{393}
$\mathrm{Imag.}\frac m{p-p^0}$ or $\mathrm{Imag.}\frac1{\xi+\omega_1i}$ has the minus sign. Therefore, for $\mathrm{Real.}\,p=0$, one obtains
\[
\varepsilon=-1\text{ if }(\sigma)=-1,\quad\varepsilon=0\text{ if }(\sigma)=+1,\quad\text{i.e. }\varepsilon=-\tfrac12+\tfrac12(\sigma)
\]
and in the same way, for $\mathrm{Imag.}\,p=0$,
\[
\begin{gathered}
\xi-\xi_1=\sigma b\mathfrak B\omega_1-\sigma\omega_4,\quad\mathrm{Imag.}\frac m{p-p^0}=\mathrm{Imag.}\frac1{\xi+\omega_1i},\\
\varepsilon=-\tfrac12+\tfrac12(\sigma)\quad\text{thus }g'=-1+(\sigma)
\end{gathered}
\]

For the fourth side, $P^0=0$, $\mathrm{Imag.}\frac M{P-P^0}=\mathrm{Imag.}\frac1{\xi+\omega_4i}$,
\[
\begin{aligned}
&\text{thus for }\mathrm{Real.}(i^\lambda P)=0;&&\xi-\xi_4=-\sigma a\mathfrak A\omega_4+\sigma\omega_1,\\
&&&\varepsilon=-\tfrac12+\tfrac12(\sigma a\mathfrak A)\\
&\text{and for }\mathrm{Imag.}(i^\lambda P)=0;&&\xi-\xi_4=-\sigma b\mathfrak B\omega_4+\sigma\omega_1,\\
&&&\varepsilon=-\tfrac12+\tfrac12(\sigma b\mathfrak B)
\end{aligned}
\]
and consequently $G_{\prime}=-1+\tfrac12(\sigma a\mathfrak A)+\tfrac12(\sigma b\mathfrak B)$, or, since $\rho=a\mathfrak A+b\mathfrak B$, $G_{\prime}=-1+\tfrac12(\rho\sigma)+\tfrac12(\rho\sigma ab\mathfrak A\mathfrak B)$.

The part $R_1^\lambda$ of $4\Theta^\lambda$ arising from the portion of the boundary infinitely near $0$ is therefore
\[
R_1^\lambda=+G_{\prime}-g'=-(\sigma)+\tfrac12(\rho\sigma)+\tfrac12(\rho\sigma ab\mathfrak A\mathfrak B)
\]

Similar considerations give, for the parts $R_2^\lambda$, $R_3^\lambda$, $R_4^\lambda$, which have the same relations to the points $\tfrac12m,\tfrac12m+\tfrac12Mi^\lambda,\tfrac12Mi^\lambda$, respectively, as $R_1^\lambda$ has to $0$, when $\lambda$ is even,
\[
\begin{aligned}
R_2^\lambda&=-g''-i^\lambda G''=-\tfrac12(b)-\tfrac12(\mathfrak B),\\
R_3^\lambda&=-i^\lambda G'+i^\lambda g_{\prime}=i^\lambda(\sigma)-\tfrac12i^\lambda(\rho\sigma)-\tfrac12i^\lambda(\rho\sigma ab\mathfrak A\mathfrak B)\\
R_4^\lambda&=+i^\lambda g_{\prime\prime}+G_{\prime\prime}=\tfrac12i^\lambda(b)+\tfrac12i^\lambda(\mathfrak B)
\end{aligned}
\]
and when $\lambda$ is odd,
\[
\begin{aligned}
R_2^\lambda&=-g''-i^{1-\lambda}H''=-\tfrac12(b)-\tfrac12(\mathfrak B),\\
R_3^\lambda&=-i^{1-\lambda}H'-i^{\lambda-1}h_{\prime}=0,\\
R_4^\lambda&=-i^{\lambda-1}h_{\prime\prime}+G_{\prime\prime}=\tfrac12i^{\lambda+1}(a)+\tfrac12i^{\lambda+1}(\mathfrak A)
\end{aligned}
\]

[V.] Art. $[14.]$ The sums of the $\varepsilon$ formed according to Rule III are evaluated by means of the equation
\[
\begin{aligned}
&\mathrm{III},\quad\Sigma\varepsilon\text{ from all }-4\Sigma\varepsilon\text{ from those where }y\text{ integer, }[x]\text{ even}\\
&=\mathrm{III},\quad\Sigma[-\mathrm{Int.}(p-m\omega)+\mathrm{Int.}(p+m\omega)]
\end{aligned}
\]
which is readily verified from the definition of the $\varepsilon$. Here $p$ takes all values satisfying the conditions stated under Rule III. Except for the two intensors corresponding to the smallest $(\xi^{*})$ and the greatest admissible value $(\xi^{**})$ of $\xi$, respectively, which equal
\[
\begin{gathered}
-\mathrm{Int.}(p^{*}-m\omega)\text{ and }+\mathrm{Int.}(p^{**}+m\omega)\text{ or}\\
-\mathrm{Int.}(\omega m)\text{ and }+\mathrm{Int.}(\tfrac12-\omega)m
\end{gathered}
\]
these intensors can always be paired as $+\mathrm{Int.}(p'+m\omega)$ and $-\mathrm{Int.}(p''-m\omega)$, so that between $\xi'$ and $\xi''$, corresponding to $p'$ and $p''$, there is no value of $\xi$ making the real or imaginary part of $p$ an integer. Thus the two intensors always cancel one another.

\SourceLeaf{394}
[V.] Art. $[15.]$ We have
\[
\begin{aligned}
\mathrm{II},\quad&+\Sigma\varepsilon\mathrm{Int.}\,ip\text{ where }Y\text{ integer }[X]\text{ even},\\
&-\Sigma\varepsilon\mathrm{Int.}\,ip\text{ where }X\text{ integer }[Y]\text{ odd}\\
={}&\Sigma[-\mathrm{Int.}\,i(p-m\omega)+\mathrm{Int.}\,i(p+m\omega)]\\
&\text{for those }p,\text{ for which }x\text{ or }y\text{ integer, }[X]\text{ even}\\
&Y\text{ odd, }0<\xi<\tfrac12,\quad\eta=\omega\\
&+\mathrm{Int.}\,i(p^{*}-m\omega)\text{ if }[X^{*}]\text{ even }[Y^{*}]\text{ odd}\\
&-\mathrm{Int.}\,i(p^{**}+m\omega)\text{ if }[X^{**}]\text{ even }[Y^{**}]\text{ odd}
\end{aligned}
\]
as one readily verifies by performing the summation on the second side of the equation, by the method used in the preceding note, over each interval from $\xi_{\prime}$ to $\xi_{\prime\prime}$ so small that between $\xi_{\prime}$ and $\xi_{\prime\prime}$ there is no $\xi$ making the real or imaginary part of the corresponding $P$ an integer. Using the $\varepsilon$ formed according to Rule III gives the second side of this equation the form stated in Article $15$.

[V.] Art. $[15.]$ The sums of the $\lambda'\varepsilon$ formed according to Rule IV are transformed into the sums of the $\varepsilon$ from V by considerations of the same kind as in the last note, also using the equation
\[
\begin{aligned}
&\mathrm V,\quad\Sigma\varepsilon\text{ from all, }-4\Sigma\varepsilon\text{ where }x\text{ integer }[y]\text{ even},\\
&=+\mathrm{Int.}(\tfrac12im+m\omega)-\mathrm{Int.}(\tfrac12im+\tfrac12-\omega m)
\end{aligned}
\]
which corresponds to the preceding evaluation of the sum of the $\varepsilon$ in Rule III.

[V.] Art. $[17.]$ Define the auxiliary quantities $U$, $T$, $L$, $V$ by the equations
\[
\begin{aligned}
U={}&1-2(B)-(AB)+(\alpha a)+(\beta b)+(\alpha b)-(\beta a)\\
&+(\alpha\beta ab)-(\alpha\beta abAB)\quad\text{or}\\
U={}&(1+(A))(1-(B))(1-(a\beta)+(b\beta)+(\alpha\beta))
\end{aligned}
\]
\[
\begin{gathered}
\text{since }(\alpha a)+(\beta b)=(A)+(A\alpha\beta ab),\\
(\alpha b)-(\beta a)=(B)-(B\alpha\beta ab)\text{},
\end{gathered}
\]
\[
\begin{aligned}
T&=-2+(a)+(b)-(\beta)-2(B)-(\alpha ab)\quad\text{or}\\
T&=-2+(a)+(b)-(\beta)-2(B)-(aA)+(bB)-(\beta AB)\\
L&=(1+(A))(1+(B))(1+(\beta))-(1-(A))(1-(B))(1-(\beta))\\
V&=(1+(A))(1-(B))(-2(a)+2(\beta)+(ab)+(\alpha\beta))\quad\text{or}\\
V&=(1+(A))(1-(B))(-(a)+(b)+(\beta)+(ab)-(\beta ab)+(\alpha\beta))
\end{aligned}
\]
\[
\text{since }(a)+(b)=(\beta)+(\beta ab)\text{ if }A\text{ positive }B\text{ negative}
\]
and denote by $W'$, $S'$, $Q'$ the quantities into which $W$, $S$, $Q$ in the expression for $\mathrm{Dec.}\frac mM$ in Article $14$ pass when $m$ and $M$, and hence $\alpha+\beta i$ and $\alpha-\beta i$, are interchanged. Then

\SourceLeaf{395}
\[
\begin{aligned}
2W'&=-5(B)-(AB)\quad\text{if }\frac{M-1}{2}\text{ even}\\
2W'&=-(B)-(AB)\quad\text{if }\frac{M-1}{2}\text{ odd}\\
2S'&=-(\alpha\beta abAB)-(\beta)-(B)+(b)\\
2Q'+2S'+2W'&=2T+U\quad\text{if }\frac{M-1}{2}\text{ even}\\
2Q'+2S'+2W'&=-4+4(a)+U\quad\text{if }\frac{M-1}{2}\text{ odd}\\
8\psi&=8(q+r+s+w)-(2Q'+2S'+2W')
\end{aligned}
\]

If one replaces $8(q+r+s+w)$ here by its value stated in Article $15$, but writes it in the form
\[
\begin{array}{lll}
2T+4L+V&\text{if }\dfrac{M-1}{2}\text{ even}&\dfrac{m-1}{2}\text{ even}\\[4pt]
2T-16+16(A)+V&\text{if }\dfrac{M-1}{2}\text{ even}&\dfrac{m-1}{2}\text{ odd}\\[4pt]
-4+4(a)+V&\text{if }\dfrac{M-1}{2}\text{ odd}&\dfrac{m-1}{2}\text{ even}\\[4pt]
-20+4(a)+V&\text{if }\dfrac{M-1}{2}\text{ odd}&\dfrac{m-1}{2}\text{ odd}
\end{array}
\]
and notes that
\[
V-U=-\tfrac12(1+(a))(1+(A))(1+(\alpha))(1-(b))(1-(B))(1-(\beta))
\]
holds, one obtains for $\psi$ the method of determination given in Article $17$.

[VI.] Art. $3$. The incomplete reference may be taken to refer to Article $4$ of Fragment III.
\begin{flushright}\textsc{Schering.}\end{flushright}
\BookRule{.3\textwidth}
